\documentclass[10pt,letterpaper]{article}
\usepackage[margin=0.15in]{geometry}
\usepackage{multicol}
\usepackage{amsmath,amssymb}
\usepackage{enumitem}
\usepackage{blindtext}
\usepackage{graphicx}

\usepackage{setspace}

\setlength{\parskip}{0em}
\setlength{\parindent}{0pt}
\setstretch{0.98}
\setlist[itemize,enumerate]{noitemsep,topsep=0pt,parsep=0pt,partopsep=0pt}
\pagestyle{empty}

\begin{document}

\begin{multicols}{2}

\textbf{Axioms of Probability}
\begin{enumerate}
    \item Non-negativity: For any event \( A \), \( P(A) \geq 0 \).
    \item Normalization: \( P(S) = 1 \), where \( S \) is the sample space.
    \item Additivity: For any two mutually exclusive events \( A \) and \( B \), \( P(A \cup B) = P(A) + P(B) \).
    This can be extended to $P(\bigcup_{i=1}^n A_i) = \sum_{i=1}^n P(A_i)$ for mutually exclusive events \( A_1, A_2, \ldots, A_n \).
\end{enumerate}

\vspace{0.1cm}
\hrule
\vspace{0.1cm}

\textbf{Basic Properties of Probability:}\\
$\forall A \in S, P(A) < 1$\\
$P(A^c) = 1 - P(A)$\\
P($\emptyset$) = 0\\
$P(A \cup B) = P(A) + P(B) - P(A \cap B)$\\
$A \subseteq B \implies P(A) \leq P(B)$

\vspace{0.1cm}
\hrule
\vspace{0.1cm}

\textbf{Counting:}\\
\textbf{Permutation} --- An arrangement of objects in a specific order. $P(n, r) = \frac{n!}{(n-r)!}$\\
\textbf{Combination} --- A selection of objects without regard to the order. $C(n, r) = \frac{n!}{r!(n-r)!}$\\
\textbf{Binomial Coefficient} --- The number of ways to choose \( k \) successes in \( n \) independent Bernoulli trials. 
$\binom{n}{k} = C(n, k) = \frac{n!}{k!(n-k)!}$\\
\textbf{Number of Subsets} --- The number of subsets of a set with \( n \) elements is \( 2^n \).

\vspace{0.1cm}
\hrule
\vspace{0.1cm}

\textbf{Conditional Probability:}\\
$P(A|B) = \frac{P(A \cap B)}{P(B)}$\\
\textbf{Multiplication Rule:} --- $P(A \cap B) = P(A|B)P(B) = P(B|A)P(A)$\\
$P(\bigcap_{i=1}^n A_i) = P(A_1)P(A_2|A_1)P(A_3|A_1 \cap A_2)...P(A_n|\bigcap_{i=1}^{n-1} A_i)$\\
\textbf{Law of Total Probability:} --- If \( B_1, B_2, \ldots, B_n \) are mutually exclusive and exhaustive events, then for any event \( A \):
$P(A) = \sum_{i=1}^n P(A|B_i)P(B_i)$\\
\textbf{Bayes' Theorem:} --- \\
$P(B_j|A) = \frac{P(A|B_j)P(B_j)}{P(A)} = \frac{P(A|B_j)P(B_j)}{\sum_{i=1}^n P(A|B_i)P(B_i)}$

\vspace{0.1cm}
\hrule
\vspace{0.1cm}

\textbf{Independence:}\\
Two events \( A \) and \( B \) are independent if and only if $P(A|B) = P(A)$ and \( P(A \cap B) = P(A)P(B) \).\\
If $A$ and $B$ are independent, then $B$ and $A$, $A$ and $B^c$, $A^c$ and $B$, $A^c$ and $B^c$ are also independent.

\vspace{0.1cm}
\hrule
\vspace{0.1cm}

\textbf{Random Variables}
A RV is a function that assigns a real number to each outcome in the sample space of a random experiment.\\
\textbf{Discrete RV} --- A random variable that can take on a countable number of distinct values.\\
\textbf{Continuous RV} --- A random variable that can take on any value within a given range or interval.\\
$X: S \rightarrow \mathbb{R}$ Where S is the event space and \( X \) is the random variable.\\
\textbf{PMF} --- $f(x) = P(X = x)$ (discrete)\\
$P(a \leq X \leq b) = \int_{a}^{b} f(x)dx$ (continuous)

\vspace{0.1cm}
\hrule
\vspace{0.1cm}

\textbf{CDF} --- $F(x) = P(X \leq x) = \sum_{y \leq x} P(X = y)$ (discrete)\\
$F(x) = P(X \leq x) = \int_{-\infty}^{x} f(t)dt$ (continuous)\\
$P(a \leq X \leq b) = F(b) - F(a)$\\
\textbf{Bernoulli Random Variable} --- A discrete RV that takes the value 1 with probability \( p \) and the value 0 with probability \( 1-p \).\\
\textbf{Expected Value} --- $E(X) = \mu_X = \sum_x xP(X=x)$ or $= \int_{-\infty}^{\infty} xf(x)dx$\\
$E(h(X)) = \sum_x h(x)P(X=x)$ or $E(h(X)) = \int_{-\infty}^{\infty} h(x)f(x)dx$\\
$E(X + Y) = E(X) + E(Y)$\\
$E(X^2) = \sum_x x^2p(x)$ or $= \int x^2f(x)dx$\\
\textbf{Variance} --- $\sigma_X^2 = E[(X - \mu_X)^2] = E(X^2) - (E(X))^2$ or $\sigma_X^2 = \int (x - \mu_X)^2 f(x) dx$\\
\textbf{Standard Deviation} --- $\sigma_X = \sqrt{Var(X)}$

% \vspace{0.1cm}
% \hrule
% \vspace{0.1cm}
\vspace{1cm}

\textbf{Binomial Distribution} --- Models the number of successes in \( n \) independent Bernoulli trials, each with success probability \( p \).\\
$X \sim Bin(n, p)$\\
\textbf{PMF:} --- $P(X = k) = b(x; n,p) = \binom{n}{k} p^k (1-p)^{n-k}$ for $ k = 0, 1, \ldots, n $ else $0$\\
\textbf{CDF:} --- $B(x; n,p) = P(X \leq x) = \sum_{y=0}^{\lfloor x \rfloor} b(y; n,p)$\\
\textbf{Mean:} --- $E(X) = np$\\
\textbf{Variance:} --- $Var(X) = np(1-p)$

\vspace{0.2cm}

\textbf{Negative Binomial Distribution} --- Models the number of trials needed to achieve \( r \) successes in a sequence of independent Bernoulli trials, 
each with success probability \( p \).\\
$X \sim NB(r, p)$\\
\textbf{PMF:} --- $P(X = k) = \binom{k-1}{r-1} p^r (1-p)^{k-r}$ for $ k = r, r+1, \ldots $ else $0$\\
\textbf{CDF:} --- $NB(x; r,p) = P(X \leq x) = \sum_{y=r}^{\lfloor x \rfloor} \binom{y-1}{r-1} p^r (1-p)^{y-r}$\\
\textbf{Mean:} --- $E(X) = \frac{r}{p}$\\
\textbf{Variance:} --- $Var(X) = \frac{r(1-p)}{p^2}$

\vspace{0.1cm}
\hrule
\vspace{0.1cm}

\textbf{Geometric Distribution} --- A special case of the negative binomial distribution where \( r = 1 \). 
It models the number of trials needed to get the first success.\\
$X \sim Geo(p)$\\
\textbf{PMF:} --- $P(X = k) = (1-p)^{k-1} p$ for $ k = 1, 2, \ldots $ else $0$\\
\textbf{CDF:} --- $G(x; p) = P(X \leq x) = 1 - (1-p)^{\lfloor x \rfloor}$\\
\textbf{Mean:} --- $E(X) = \frac{1}{p}$\\
\textbf{Variance:} --- $Var(X) = \frac{1-p}{p^2}$

\vspace{0.1cm}
\hrule
\vspace{0.1cm}

\textbf{Hypergeometric Distribution} --- Models the number of successes in \( n \) draws from a finite population of size \( N \) containing \( K \) successes,
without replacement.\\
$X \sim Hyp(N, K, n)$\\
\textbf{PMF:} --- $h(x; n, N, K) = P(X = k) = \frac{\binom{K}{k} \binom{N-K}{n-k}}{\binom{N}{n}}$ for $ \max(0, n-(N-K)) \leq k \leq \min(n, K) $ else $0$\\
\textbf{CDF:} --- $H(x; n, N, K) = P(X \leq x) = \sum_{y=\max(0, n-(N-K))}^{\lfloor x \rfloor} \frac{\binom{K}{y} \binom{N-K}{n-y}}{\binom{N}{n}}$\\
\textbf{Mean:} --- $E(X) = n \frac{K}{N}$\\
\textbf{Variance:} --- $Var(X) = n \frac{K}{N} \frac{N-K}{N} \frac{N-n}{N-1}$

\vspace{0.1cm}
\hrule
\vspace{0.1cm}

\textbf{Poisson Distribution} --- Models the number of events occurring in a fixed interval of time or space, given the average rate \( \lambda \) of occurrence.\\
Used to model the number of times a rare event occurs in a large population.\\
$X \sim Pois(\lambda)$\\
\textbf{PMF:} --- $P(X = k) = \frac{\lambda^k e^{-\lambda}}{k!}$ for $ k = 0, 1, 2, \ldots $ else $0$\\
\textbf{CDF:} --- $P(X \leq x) = \sum_{y=0}^{\lfloor x \rfloor} \frac{\lambda^y e^{-\lambda}}{y!}$\\
\textbf{Mean:} --- $E(X) = \lambda$\\
\textbf{Variance:} --- $Var(X) = \lambda$ 

\vspace{0.1cm}
\hrule
\vspace{0.1cm}

\textbf{Normal} --- $X \sim \mathcal{N}(\mu, \sigma^2)$\\
\textbf{PDF:} --- $f(x) = \frac{1}{\sigma \sqrt{2 \pi}} e^{-\frac{(x - \mu)^2}{2\sigma^2}}$ for $ -\infty < x < \infty $ else $0$\\
\textbf{Standard Normal} --- $Z \sim \mathcal{N}(0, 1) \qquad Z = \frac{X - \mu}{\sigma}$\\
$P(a \leq X \leq b) = P(\frac{a - \mu}{\sigma} \leq Z \leq \frac{b - \mu}{\sigma}) = \Phi(\frac{b - \mu}{\sigma}) - \Phi(\frac{a - \mu}{\sigma})$\\
$z_\alpha$ is the z-score such that $P(Z \geq z_\alpha) = \alpha$ or the $(1 - \alpha) \cdot 100$th percentile

\vspace{0.1cm}
\hrule
\vspace{0.1cm}

\textbf{Exponential} --- $X \sim Exp(\lambda)$\\
\textbf{PDF:} --- $f(x) = \lambda e^{-\lambda x}$ for $ x \geq 0 $ else $0$\\
\textbf{CDF:} --- $F(x) = 1 - e^{-\lambda x}$ for $ x \geq 0 $ else $0$\\
\textbf{Mean:} --- $E(X) = \frac{1}{\lambda}$\\
\textbf{Variance:} --- $Var(X) = \frac{1}{\lambda^2}$

% \vspace{0.1cm}
% \hrule
% \vspace{0.1cm}
\vspace{1.5cm}

\textbf{Uniform} --- $X \sim U(a, b)$\\
\textbf{PDF:} --- $f(x) = \frac{1}{b - a}$ for $ a \leq x \leq b $ else $0$\\
\textbf{CDF:} --- $F(x) = \frac{x - a}{b - a}$ for $ a \leq x \leq b $, $0$ for $ x < a $, $1$ for $ x > b $\\
\textbf{Survival Function:} --- $S(x) = 1 - F(x) = \frac{b - x}{b - a}$ for $ a \leq x \leq b $, $1$ for $ x < a $, $0$ for $ x > b $\\
\textbf{Mean:} --- $E(X) = \frac{a + b}{2}$\\
\textbf{Variance:} --- $\sigma^2 = \frac{(b - a)^2}{12}$

\vspace{0.1cm}
\hrule
\vspace{0.1cm}

\textbf{Jointly Distributed Random Variables:}\\
$P(x, y) = P(X = x, Y = y)$\\
Trick: $P(X=Y) + P(Y > X) + P(Y < X) = 1$\\
\textbf{Marginal PMF} --- $P_X(x) = \sum_y P(x, y)$ and $P_Y(y) = \sum_x P(x, y)$\\
$X$ and $Y$ are independent $\iff P(x, y) = P_X(x)P_Y(y)$\\
\textbf{Conditional PMF} --- $P_{Y|X}(y|x) = \frac{P(x, y)}{P_X(x)}$\\
\textbf{Expected Value} --- $E(h(X,Y)) = \sum_x \sum_y h(x,y) P(x,y)$\\
$E(h(X,Y)) = \int \int h(x,y) f(x,y) dx dy$\\
\textbf{Covariance} --- $Cov(X,Y) = E[(X - \mu_X)(Y - \mu_Y)] = E(XY) - E(X)E(Y)$\\
$Cov(X,Y) = 0$ if \( X \) and \( Y \) are independent.\\
\textbf{Correlation Coefficient} --- $\rho_{XY} = \frac{Cov(X,Y)}{\sigma_X \sigma_Y}$\\
$-1 \leq \rho_{XY} \leq 1$\\
\textbf{Properties} $Cov(X,X) = Var(X)$ $Cov(X,Y) = Cov(Y,X)$ $Cov(aX,bY) = abCov(X,Y)$
$Cov(X + Y, Z) = Cov(X,Z) + Cov(Y,Z)$ $\exists a,b \in \mathbb{R} \ni Y = aX + b$

\vspace{0.1cm}
\hrule
\vspace{0.1cm}

\textbf{Random Samples and Statistics:}\\
A random sample of size \( n \) is a set of \( n \) independent and identically distributed (i.i.d.) random variables \( X_1, X_2, \ldots, X_n \) drawn from a population with a common distribution.\\
\textbf{Sample Mean} --- $\bar{X} = \frac{1}{n} \sum_{i=1}^n X_i$\\
\textbf{Sample Variance} --- $S^2 = \frac{1}{n-1} \sum_{i=1}^n (X_i - \bar{X})^2$\\
$E(\bar{X}) = \mu$\\
$Var(\bar{X}) = \frac{\sigma^2}{n}$\\
When sampling from a normal population, $\bar{X} \sim \mathcal{N}(\mu, \frac{\sigma^2}{n})$

\vspace{0.1cm}
\hrule
\vspace{0.1cm}

\textbf{Linear Combinations} --- $Y = a_1X_1 + a_2X_2 + ... + a_nX_n$\\
$E(Y) = a_1E(X_1) + a_2E(X_2) + ... + a_nE(X_n)$\\
$Var(Y) = a_1^2Var(X_1) + a_2^2Var(X_2) + ... + a_n^2Var(X_n)$ if \( X_i \) are independent.\\
If \( X_i \sim \mathcal{N}(\mu_i, \sigma_i^2) \), then $Y \sim \mathcal{N}(a_1\mu_1 + a_2\mu_2 + ... + a_n\mu_n, a_1^2\sigma_1^2 + a_2^2\sigma_2^2 + ... + a_n^2\sigma_n^2)$

\vspace{0.1cm}
\hrule
\vspace{0.1cm}

\textbf{Central Limit Theorem:}\\
$n > 30 \implies \frac{\bar{X} - \mu}{\sigma / \sqrt{n}} \sim \mathcal{N}(0, 1)$\\
\textbf{Point Estimation:}\\
Sample Mean: $\bar{X} = \frac{1}{n} \sum_{i=1}^n X_i$\\
Sample Proportion: $\hat{p} = \frac{M}{n}$, M = \# of successes\\
Sample Variance: $S^2 = \frac{1}{n-1} \sum_{i=1}^n (X_i - \bar{X})^2$\\
\textbf{Quantitative} ($\mu$):\\
Point Estimator: $\hat{\theta} = \bar{X}$\\
Bias: $E(\hat{\theta}) - \theta$ Bias of 0 means unbiased estimator\\
Standard Error: $\sigma_{\hat{X}} = \frac{\sigma}{\sqrt{n}}$\\
Est Standard Error: $\frac{S}{\sqrt{n}}$\\
$\text{Var}(\hat{\theta}) = \frac{\theta(1-\theta)}{n}$\\
$P(|\bar{X} - \mu| \leq c) = P(-c \leq \bar{X} - \mu \leq c) = P\left(-\frac{c}{\sigma / \sqrt{n}} \leq Z \leq \frac{c}{\sigma / \sqrt{n}}\right)$\\
$= \Phi_Z\left(\frac{c}{\sigma / \sqrt{n}}\right) - \Phi_Z\left(-\frac{c}{\sigma / \sqrt{n}}\right) = 2\Phi_Z\left(\frac{c}{\sigma / \sqrt{n}}\right) - 1$\\
\textbf{Qualitative} (p):\\
Point Estimator: $\hat{p} = \frac{M}{n}$\\
Bias: $E(\hat{p}) - p$ Bias of 0 means unbiased estimator\\
Standard Error: $\sqrt{\frac{p(1-p)}{n}}$\\
Est Standard Error: $\sqrt{\frac{\hat{p}(1-\hat{p})}{n}}$

% \vspace{0.1cm}
% \hrule
% \vspace{0.1cm}
\vspace{1cm}

\textbf{Confidence Intervals:}\\
Confidence Level: $100(1-\alpha)\%$\\
$CI_\mu = (\bar{X} \pm z_{\frac{\alpha}{2}} \cdot \frac{\sigma}{\sqrt{n}})$, $z_{\frac{\alpha}{2}}$ is the z-score such that $\Phi_Z\left(z_{\frac{\alpha}{2}}\right) = 1 - \frac{\alpha}{2}$\\
$99\% \rightarrow \alpha = 0.01 \rightarrow z_{\frac{\alpha}{2}} = 2.576$\\
$95\% \rightarrow \alpha = 0.05 \rightarrow z_{\frac{\alpha}{2}} = 1.96$\\
$90\% \rightarrow \alpha = 0.10 \rightarrow z_{\frac{\alpha}{2}} = 1.645$\\      
Higher confidence levels lead to wider confidence intervals. Larger sample sizes lead to 
narrower confidence intervals. Larger standard deviations lead to wider confidence intervals.\\
\textbf{Confidence Intervals for Large Samples:}\\
If $n \geq ~40$ then we can replace $\sigma$ with $S$ in the confidence interval formula in cases where $\sigma$ is unknown.\\
$Z = \frac{\bar{X} - \mu}{S / \sqrt{n}} \sim \mathcal{N}(0, 1)$ approximately by CLT.\\
$CI_\mu = (\bar{X} \pm z_{\frac{\alpha}{2}} \cdot \frac{S}{\sqrt{n}})$\\
$CI_p = \left(\hat{p} \pm z_{\frac{\alpha}{2}} \cdot \sqrt{\frac{\hat{p}(1-\hat{p})}{n}}\right)$

\vspace{0.1cm}
\hrule
\vspace{0.1cm}

\textbf{Hypothesis Testing:}\\
\textbf{Statistical Hypothesis} --- A claim about a population parameter.\\
\textbf{Test Statistic} --- A statistic calculated from sample data that is used to make decisions about the hypothesis.\\
\textbf{Null Hypothesis} ($H_0$) --- The hypothesis that there is no effect or difference, and it is the hypothesis that is tested.\\
\textbf{Alternative Hypothesis} ($H_a$)--- The hypothesis that there is an effect or difference.\\
\textbf{Conclusion} --- Reject $H_0$ or Fail to Reject $H_0$.\\
\textbf{Type I Error} --- Rejecting the null hypothesis when it is true.\\
\textbf{Type II Error} --- Failing to reject the null hypothesis when it is false.\\
\textbf{Significance Level} ($\alpha$) --- The probability of making a Type I error. Common choices are 0.05, 0.01, and 0.10.\\
\textbf{Rejection Region} --- The set of values for the test statistic that leads to the rejection of the null hypothesis. $\text{p-value} \leq \alpha \implies \text{Reject } H_0$

\vspace{0.1cm}
\hrule
\vspace{0.1cm}

\textbf{P-Values}\\
The probability of obtaining a test statistic as extreme as the one observed, assuming the null hypothesis is true.\\
\textbf{Left-Tailed Test} ($H_a: \mu < \mu_0$):\\
$p\text{-value} = P(Z \leq z_{obs}) = \Phi_Z(z_{obs})$\\
\textbf{Right-Tailed Test} ($H_a: \mu > \mu_0$):\\
$p\text{-value} = P(Z \geq z_{obs}) = 1 - \Phi_Z(z_{obs})$\\
\textbf{Two-Tailed Test} ($H_a: \mu \neq \mu_0$):\\
$p\text{-value} = 2P(Z \geq |z_{obs}|) = 2(1 - \Phi_Z(|z_{obs}|))$

\vspace{0.1cm}
\hrule
\vspace{0.1cm}

\textbf{Tests for a Population Mean:}\\
\textbf{Case 1:}
Sampling from a Normal population with known $\sigma$ and unknown $\mu$. Suppose we want to test $H_0 : \mu = \mu_0$. In this case our test
statistic is $Z = \frac{\bar{X} - \mu_0}{\sigma / \sqrt{n}}$. $\bar{X}$ is the sample mean and $\mu_0$ is the hypothesized population mean under $H_0$.\\
\textbf{Case 2:}
If $\sigma$ is unknown and $n \geq 40$, we can use the same test statistic as in Case 1, but replace $\sigma$ with $S$ (the sample standard deviation).
$Z = \frac{\bar{X} - \mu_0}{S / \sqrt{n}}$ approximately follows a standard normal distribution by the CLT.

\vspace{0.1cm}
\hrule
\vspace{0.1cm}

\textbf{Tests for a Population Proportion:}\\
We want to test $H_0 : p = p_0$. Our test statistic is $Z = \frac{\hat{p} - p_0}{\sqrt{p_0(1-p_0)/n}}$. $\hat{p}$ is the sample proportion and $p_0$ is the hypothesized population proportion under $H_0$.\\

\vspace{9999pt}

\end{multicols}

\end{document}